El presente artículo examina la relevancia y la interacción de cinco componentes fundamentales en el desarrollo de software contemporáneo: Git y GitHub, la Guía de Scrum 2020, los patrones de diseño, el patrón Modelo-Vista-Controlador (MVC) y el lenguaje de programación Java. Se adopta un enfoque descriptivo y analítico para mostrar cómo estas herramientas, metodologías y arquitecturas favorecen la colaboración, la trazabilidad y la calidad del código en proyectos de diversa escala. A lo largo del estudio se destaca que Git, combinado con la plataforma GitHub, asegura un control de versiones eficiente y un trabajo en equipo distribuido. Asimismo, la aplicación de Scrum según la guía 2020 promueve iteraciones cortas, transparencia y adaptabilidad a los cambios. Por su parte, los patrones de diseño y el patrón MVC refuerzan la modularidad, la reutilización de componentes y la separación de responsabilidades, mientras que Java aporta portabilidad, robustez y un ecosistema amplio que facilita la construcción de soluciones escalables. Los resultados de este análisis demuestran que la integración coordinada de estas prácticas constituye una base sólida para proyectos académicos y profesionales, incrementando la eficiencia, la mantenibilidad y la capacidad de respuesta frente a las demandas del entorno tecnológico actual.